\textbf{\underline{\large{4.2: The Mean Value and Rolle's Theorem}}} \par

The Mean Value Theorem is an interesting piece of the history of calculus that was used to prove a lot of what we take for granted. The Mean Value THeorem was used to prove that a derivative being positive or negative told you that the function was increasing or decreasing, respectively. Of course, this led directly to the First Derivative Test and the intervals of concavity. \par

\begin{center}
    \fbox{\fbox{\begin{minipage}{0.96 \textwidth}
        \vspace{11pt}
        \begin{center}
            \textbf{Mean Value Theorem}
        \end{center}
        \vspace{11pt}
        If $f$ is a function that satisfies these two hypotheses:
        \vspace{11pt}
        \begin{center}
            \begin{tabular}{l}
                1) $f$ is continuous on the closed interval $[a, \, b]$. \\[11pt]
                2) $f$ is differentiable on the closed interval $(a, \, b)$.
            \end{tabular}
        \end{center}
        \vspace{11pt}
        Then there exists a number $c$ in the interval $(a, \, b)$ such that \begin{align*}
            f'(c) = \dfrac{f(b) - f(a)}{b - a}.
        \end{align*}
    \end{minipage}}}
\end{center}

Again, translating from math to english, this simply states that if you have a smooth, continuous curve, the slope of the line connecting the endpoints has to equal the slope of a tangent somewhere in that interval. Alternatively, it says that the secant line (a line through two distinct points on a curve) through the endpoints will have the same slope as a tangent line at some point on the curve. \\

\begin{center}
    \includegraphics[width = 0.5\textwidth]{\graphicsdir Chapter 4 Graphics/4.2-Graphic1.png}
\end{center}

\begin{tcolorbox}[example]
    \textbf{Ex 4.2.1: } Show that the function $f(x) = x^3 - 4x^2 + 1$ with $x \in [-1, \, 3]$ satisfies all the conditions of the Mean Value Theorem and find $c$.
\end{tcolorbox}
\begin{tcolorbox}[solution]
    \textbf{Sol 4.2.1: } Polynomials are continuous throughout their domain, so the first hypothesis is satisfied. Polynomials are also differentiable throughout their domain, so the second condition is satisfied. \begin{align*}
        f'(c) &= \dfrac{f(b) - f(a)}{b - a} \\[11pt]
        f'(c) = 3c^2 - 8c &= \dfrac{f(3) - f(-1)}{3 - (-1)} \\[11pt]
        3c^2 - 8c &= \dfrac{-8 - (-4)}{4} \\[11pt]
        3c^2 - 8c &= -1 \implies 3c^2 - 8c + 1 = 0 \\[11pt]
        c &= \boxed{2.535, \, 0.131}
    \end{align*}
    \begin{tcolorbox}[interesting]
        For this problem, we actually find \textit{two} values of $c$. This is definitely possible, as the Mean Value Theorem doesn't state how many numbers $c$ that satisfy the theorem; simply that there exists at least one number $c$.
    \end{tcolorbox}
\end{tcolorbox} \vspace{11pt}

\begin{tcolorbox}[example]
    \textbf{Ex 4.2.2: } Show that the function $f(x) = x^2 - 4x + 1$ with $x \in [0, \, 4]$ satisfies all the conditions of the Mean Value Theorem and find $c$.
\end{tcolorbox}
\begin{tcolorbox}[solution]
    \textbf{Sol 4.2.2: } Polynomials are continuous throughout their domain, so the first hypothesis is satisfied. Polynomials are also differentiable throughout their domain, so the second condition is satisfied. \begin{align*}
        f'(c) &= \dfrac{f(b) - f(a)}{b - a} \\[11pt]
        f'(c) = 2c - 4 &= \dfrac{f(4) - f(0)}{4 - 0} \\[11pt]
        2c - 4 &= \dfrac{1 - 1}{4} \\[11pt]
        2c - 4 &= 0 \\[11pt]
        c &= \boxed{2}
    \end{align*}
\end{tcolorbox}

Now, let's talk a little about Rolle's Theorem. Rolle's Theorem is a specific case of the Mean Value Theorem (though Joseph-Louis Lagrange used it to prove the Mean Value Theorem). Rolle's Theorem is extremely powerful, as it was used to prove all of the rules we have used to interpret derivatives. It used to be a very useful theorem, but is now something of a historical curiosity. \par

\begin{center}
    \fbox{\fbox{\begin{minipage}{0.96 \textwidth}
        \vspace{11pt}
        \begin{center}
            \textbf{Rolle's Theorem}
        \end{center}
        \vspace{11pt}
        If $f$ is a function that satisfies these three hypotheses:
        \vspace{11pt}
        \begin{center}
            \begin{tabular}{l}
                1) $f$ is continuous on the closed interval $[a, \, b]$. \\[11pt]
                2) $f$ is differentiable on the closed interval $(a, \, b)$. \\[11pt]
                3) $f(a) = f(b)$
            \end{tabular}
        \end{center}
        \vspace{11pt}
        Then there exists a number $c$ in the interval $(a, \, b)$ such that $f'(c) = 0$.
    \end{minipage}}}
\end{center}

Although this theorem may seem confusing at first, it is simply stating that if you have a continuous and smooth curve with an initial point and ending point at the same height, then there exists some point in the curve that has a derivative of zero. This can be simply proved by utilizing the Mean Value Theorem and noticing that $f(b) - f(a) = 0$. \begin{align*}
    f'(c) &= \dfrac{f(b) - f(a)}{b - a} \rightarrow f'(c) = \dfrac{0}{b - a} \implies f'(c) = 0
\end{align*}

Another theorem that is important to note and often confused with Rolle's Theorem or the Mean Value Theorem is the \textbf{Extreme Value Theorem}. Recall what the Extreme Value Theorem states: \par

\begin{center}
    \fbox{\fbox{\begin{minipage}{0.96 \textwidth}
        \vspace{11pt}
        \begin{center}
            \textbf{Extreme Value Theorem}
        \end{center}
        \vspace{11pt}
        If $f$ is continuous on the closed interval $[a, \, b]$, then $f(x)$ attains a maximum and a minimum at least once each.
    \end{minipage}}}
\end{center}

The final theorem that may be confused with the Mean Value Theorem is the \textbf{Intermediate Value Theorem}. Recall what the Intermediate Value Theorem states: \par

\begin{center}
    \fbox{\fbox{\begin{minipage}{0.96 \textwidth}
        \vspace{11pt}
        \begin{center}
            \textbf{Intermediate Value Theorem}
        \end{center}
        \vspace{11pt}
        If $f$ is continuous on the closed interval $[a, \, b]$, then $f(x)$ attains every height (or value) between $f(a)$ and $f(b)$.
    \end{minipage}}}
\end{center}

Now that we've gone over the Mean Value Theorem, we need to consider when this theorem might not apply to a problem --- that is, when is a function not continuous or not differentiable. We first saw the concept of non-continuous functions in precalculus, but let's formalize these discontinuities within the following table. \par

\begin{align*}
    \def\arraystretch{1.5}
    \begin{array}{|c|c|c|}
        \hline
        & \parbox{0.45\textwidth}{\centering
            \vspace{5.5pt}
            \textbf{Removable Discontinuity} \\[5.5pt]
            \Big( $\lim_{x \to a} f(x)$ Exists \Big) \\
            \vspace{5.5pt}
        }
        & \parbox{0.45\textwidth}{\centering
            \vspace{5.5pt}
            \textbf{Essential Discontinuity} \\[5.5pt]
            \Big( $\lim_{x \to a} f(x)$ Does Not Exist \Big) \\
            \vspace{5.5pt}
        } \\
        \hline
        \hspace{5.5pt}\raisebox{-.5\height}{\makebox[0pt][c]{\rotatebox{90}{\textbf{$f(a)$ Exists}}}}\hspace{5.5pt} &
        \parbox{0.45\textwidth}{\centering
            \vspace{5.5pt}
            \includegraphics[width=0.4\textwidth]{\graphicsdir Chapter 4 Graphics/4.2-Graphic4.png} \\[5.5pt]
            \textit{Point of Displacement (POD)}
            \vspace{5.5pt}
        } &
        \parbox{0.45\textwidth}{\centering
            \vspace{5.5pt}
            \includegraphics[width=0.4\textwidth]{\graphicsdir Chapter 4 Graphics/4.2-Graphic5.png} \\[5.5pt]
            \textit{Jump Discontinuity}
            \vspace{5.5pt}
        } \\
        \hline
        \hspace{5.5pt}\raisebox{-.5\height}{\makebox[0pt][c]{\rotatebox{90}{\textbf{$f(a)$ Does Not Exist}}}}\hspace{5.5pt} &
        \parbox{0.45\textwidth}{\centering
            \vspace{5.5pt}
            \includegraphics[width=0.4\textwidth]{\graphicsdir Chapter 4 Graphics/4.2-Graphic3.png} \\[5.5pt]
            \textit{Point of Exclusion (POE)}
            \vspace{5.5pt}
        } &
        \parbox{0.45\textwidth}{\centering
            \vspace{5.5pt}
            \includegraphics[width=0.4\textwidth]{\graphicsdir Chapter 4 Graphics/4.2-Graphic6.png} \\[5.5pt]
            \textit{Vertical Asymptote}
            \vspace{5.5pt}
        } \\
        \hline
    \end{array}
\end{align*}

